%! The Rational Parent Function \& Transformations — Unit 4, Lesson 4.4 — Lesson Plan
\documentclass[10pt]{article}
\usepackage{algebra2-article}
\usepackage{algebra2-boxes}
\usepackage{graphicx}   % -article does NOT load graphicx; needed for thumbnails/screenshots
\graphicspath{{images/}}
\newcommand{\TallMath}[1]{$\displaystyle #1\rule[-1.4em]{0pt}{3.2em}$}
\newcommand{\CourseName}{Algebra 2}
\newcommand{\MeetingLength}{60 minutes}
\newcommand{\UnitNumberName}{Unit 4: Rational Functions \quad}
\newcommand{\LessonNumberName}{Lesson 4.4: The Rational Parent Function \& Transformations}

\begin{document}
\begin{center}
    {\Large\bfseries \CourseName \\
    {\small\bfseries \UnitNumberName \LessonNumberName}}
\end{center}

\begin{tcolorbox}[colback=forestbg, colframe=forest, boxrule=0.9pt, arc=2mm,
    left=2mm, right=2mm, top=1.5mm, bottom=1.5mm]
    \textbf{Primary Objective:} Students can graph and describe the \textbf{rational parent function}
    $f(x)=\frac1x$ --- the two-branch \textbf{hyperbola} with a vertical asymptote $x=0$, a horizontal
    asymptote $y=0$, origin symmetry, and \emph{no intercepts at all} --- and can then move it. Given
    $g(x)=\frac{a}{x-h}+k$ they place the \textbf{asymptotes first} ($x=h$ from the denominator, which
    does the opposite of what it says; $y=k$ from the number added outside), locate the \textbf{center}
    $(h,k)$ where the two cross, and read off domain, range, and both intercepts --- knowing $a$ moves
    neither asymptote but decides how far the branches sit from the center and which pair of regions
    they occupy. Running it backwards, they \textbf{write the equation from a graph}: the asymptotes
    give $h$ and $k$, one marked point gives $a$, and a second point checks it. Finally they recognize
    a linear-over-linear function as a \emph{disguised} $\frac1x$ by rewriting its numerator in terms
    of its denominator.
    \\[2pt]
    \textbf{Standards (2023 VA SOL):} \textbf{A2.F.1a} --- distinguish the graphs of the parent
    functions (today: the rational parent among the parents already met). \textbf{A2.F.1b} --- write
    the equation of a rational function given a graph, using transformations of the parent.
    \textbf{A2.F.1c} --- graph a rational function given the equation, using $f(x)+k$, $f(kx)$,
    $f(x+k)$, and $kf(x)$. \textbf{A2.F.1e} --- compare and contrast the graphs, tables, and equations
    of the function families. Applied throughout: \textbf{A2.F.2h} (equations of vertical and
    horizontal asymptotes, from Lesson 4.0), \textbf{A2.F.2a} (domain, range, intercepts, including
    graphs with discontinuities), \textbf{A2.F.2c} (increasing/decreasing intervals), \textbf{A2.F.2g}
    (end behavior), and \textbf{A2.F.2f} (interpreting $f(x)$ in context --- the dilution and
    free-throw models). Builds on Units 1--2 transformations and Lesson 4.3's combining; it is the
    prerequisite for \textbf{4.5} (where the asymptotes stop being visible in the equation) and quietly
    seeds \textbf{4.7} (Homework 5's constant product is an inverse variation).
\end{tcolorbox}

\renewcommand{\arraystretch}{1.4}
\begin{skillbox}[Priority Ideas \& Skills]{goldbox}
\begin{tabularx}{\linewidth}{@{}X|X@{}}
\textbf{Priority Ideas \& Skills} & \textbf{Key Understandings (the ``why'')} \\[2pt]
\hline
\vspace{-6pt}
\begin{itemize}[leftmargin=1.1em, nosep, after=\vspace{2pt}]
  \item Graph and describe the parent $f(x)=\frac1x$: two branches, both asymptotes, domain, range, symmetry, end behavior.
  \item Read $a$, $h$, $k$ out of $g(x)=\frac{a}{x-h}+k$ and place the two asymptotes \emph{before} anything else.
  \item Give domain ($x\neq h$), range ($y\neq k$), and both intercepts \emph{by computing}, not by reading.
  \item Describe the effect of $a$: stretch, compression, reflection --- and no asymptote movement.
  \item Meet $f(kx)$ honestly: on this parent it collapses into $kf(x)$, because $\frac{1}{kx}=\frac{1/k}{x}$.
  \item \textbf{Write the equation from a graph} (A2.F.1b): asymptotes $\rightarrow h,k$; one point $\rightarrow a$; second point $\rightarrow$ check.
  \item Rewrite a linear-over-linear function as $\frac{a}{x-h}+k$ to expose the hidden parent.
\end{itemize}
&
\vspace{-6pt}
\begin{itemize}[leftmargin=1.1em, nosep, after=\vspace{2pt}]
  \item \textbf{The asymptotes are the parent's axes, and they travel with it.} Draw the two dashed lines and the branches have nowhere else to go.
  \item An \emph{input} shift moves a wall; an \emph{output} shift moves a floor. That is the whole lesson in one sentence.
  \item The parent has \emph{no} intercepts --- the first parent of the course for which that is true, and the fact students refuse to believe.
  \item ``Inside is backwards'' is not a new rule; it is Unit 1's rule, and it costs more here because it puts the wall on the wrong side of the axis.
  \item The horizontal asymptote is a \textbf{range} exclusion. Students restrict domains by habit and forget ranges entirely.
  \item A transformed parent \emph{never} crosses its horizontal asymptote --- stronger than 4.0's ``may cross,'' and for a reason worth stating: the numerator is a nonzero constant.
  \item $k$ \emph{is} the ratio of the leading coefficients. The degree rule from 4.0 and the general form are the same fact.
\end{itemize}
\end{tabularx}
\end{skillbox}

\begin{skillbox}[{Vocabulary, Concepts, \& Theorems}]{sky}
\renewcommand{\arraystretch}{1.3}
\begin{tabularx}{\linewidth}{@{}l X@{}}
\textbf{Rational parent function} & \TallMath{f(x)=\frac1x} --- the reciprocal. Vertical asymptote $x=0$, horizontal asymptote $y=0$, domain and range both ``all reals except $0$,'' \emph{no} intercepts. \\
\textbf{Hyperbola / branch} & The two-piece curve that is the graph of $\frac1x$; each separate piece is a \emph{branch}. \\
\textbf{General form} & \TallMath{g(x)=\frac{a}{x-h}+k} --- every transformation of the parent at once. \\
\textbf{Center} & $(h,k)$, the point where the two asymptotes cross. The parent's origin, relocated --- and \emph{not} a point of the graph. \\
\textbf{The four transformations} & $f(x)+k$ (moves the HA), $f(x+k)$ (moves the VA), $kf(x)$ (stretch/reflect), $f(kx)$ (which here equals a $kf(x)$). \\
\textbf{Origin symmetry (3.0)} & $f(-x)=-f(x)$: rotating $180^\circ$ about the origin returns the graph. $\frac1x$ is an \emph{odd} function. \\
\textbf{Degree comparison (4.0)} & Equal degrees $\Rightarrow$ HA is the ratio of the leading coefficients --- which is exactly the $k$ left over after the rewrite. \\
\end{tabularx}
\end{skillbox}

\begin{fixedskillbox}[Activate Prior Knowledge \& Spiral Review]{forestbg}
    The Warm-Up rehearses three skills, one per leg of the lesson. \textbf{(1)} \emph{The reciprocal
    table}: $y=\frac1x$ evaluated at $-4,-1,-\frac12,\frac12,1,4$ plus the undefined column at $0$,
    then domain, range, and both asymptote equations --- the parent met numerically before it is met as
    a picture. Debrief this one slowly: the answers ``no $x$-intercept'' and ``no $y$-intercept'' are
    permanent, and true of no other parent students have graphed. \textbf{(2)} \emph{The shift rules
    from Units 1--2}: the vertices of $|x-3|+1$ and $(x+2)^2-5$, then the general statement in words
    --- outside is vertical and honest, inside is horizontal and backwards. Insist on the words; today
    that rule moves a \emph{wall}. \textbf{(3)} \emph{One function, two costumes}: combine
    $2+\frac{3}{x-1}$ into $\frac{2x+1}{x-1}$ (Lesson 4.3), then find its asymptotes the Lesson 4.0 way
    --- $x=1$ and $y=2$, which are exactly the two numbers it started with. \textbf{Say nothing about
    that coincidence.} It is Section 6 of the notes, and the recognition is worth more than the hint.
\end{fixedskillbox}

\begin{skillbox}[Hook]{forestbg}
    Put two tables side by side --- $y=\frac1x$ at $x=-2,-1,0,1,2$ and $y=\frac{1}{x-3}$ at
    $x=1,2,3,4,5$ --- and have the class fill both. The output columns come out \emph{identical}:
    $-\frac12,-1,\text{undef.},1,\frac12$ down each one. \textbf{``Compare the two output columns. What
    do you notice?''} Nothing about the outputs changed; they simply arrived three inputs later. So the
    blow-up moved from $x=0$ to $x=3$, and the settling-down value stayed at $y=0$. Land the sentence
    that runs the period before you draw anything: \textbf{shifting a rational function carries its
    asymptotes with it --- the wall moved, the floor did not.} The point of doing this numerically is
    that the shift becomes a fact about the table rather than a rule to be trusted, and the ``inside is
    backwards'' direction is something students \emph{observe} ($x-3$ moved things right) instead of
    something they try to remember.
\end{skillbox}

\begin{skillbox}[Lesson]{forestbg}
\begin{multicols}{2}
\textbf{1. The parent.} Graph $y=\frac1x$ from the Warm-Up table and read it completely: domain and
range both ``all reals except $0$,'' VA $x=0$, HA $y=0$, end behavior $y\to0$ both ways. Then spend real
time on the intercepts, because the answer is \emph{none} and no previous parent behaved that way. Ask
it as two separate questions: \textbf{``when is a fraction equal to zero?''} (only when its numerator
is --- and this numerator is $1$) and \textbf{``what is $f(0)$?''} (nothing; $0$ is not in the domain).
Close with two callbacks: the graph decreases on $(-\infty,0)$ \emph{and} on $(0,\infty)$ --- two
intervals, never one, because an interval may not span a wall (4.0) --- and it has origin symmetry, so
it is odd (3.0).

\textbf{2. The two shifts.} This is the heart. $y=\frac{1}{x-3}$ moves the parent right $3$, and its
\emph{vertical} asymptote goes with it to $x=3$; $y=\frac1x+2$ moves the parent up $2$, and its
\emph{horizontal} asymptote goes to $y=2$. Draw both with the parent dashed behind. Then state the box:
\textbf{the asymptotes are the parent's axes, and they travel with it.} Every graphing question today
is answered by drawing those two dashed lines first, and students who begin by plotting points instead
will fight the vertical asymptote all period.

\textbf{3. The stretch, the flip, and the surprise.} Show $y=\frac1x$, $y=\frac3x$, and $y=-\frac1x$ on
three grids. $|a|>1$ pushes the branches away from the asymptotes, $|a|<1$ pulls them in, $a<0$ sends
them to the opposite pair of regions --- and \emph{no asymptote moves}. Then take two minutes on the
fourth transformation, which A2.F.1c names and which behaves strangely here:
$f(4x)=\frac{1}{4x}=\frac{1/4}{x}$, so a horizontal compression by $4$ \emph{is} a vertical compression
by $\frac14$. On this parent --- and only on this parent --- $f(kx)$ and $kf(x)$ give the same family
of curves. That is the honest answer to ``why is there only one letter for stretching?''

\columnbreak

\textbf{4. The general form.} Assemble $g(x)=\frac{a}{x-h}+k$ and read it in a fixed order: $h$ gives
the VA, $k$ gives the HA, $(h,k)$ is the \textbf{center} where they cross (and is not on the graph),
$a$ does the rest. Work the anchor $g(x)=\frac{2}{x-3}+1$ completely: VA $x=3$, HA $y=1$, center
$(3,1)$, domain $x\neq3$, range $y\neq1$, $y$-intercept $\frac13$, $x$-intercept $(1,0)$ --- and insist
the intercepts are \emph{computed}. Then the question that separates today from 4.0: \textbf{``does
this graph ever cross its horizontal asymptote?''} Setting $g(x)=1$ requires $\frac{2}{x-3}=0$, which is
impossible. A transformed parent \emph{never} crosses its HA --- stronger than 4.0's ``may cross,'' and
for a nameable reason.

\textbf{5. Backwards: the equation from a graph (A2.F.1b).} Read the asymptotes to get $h$ and $k$ for
free, substitute one clearly marked point to solve for $a$, then \emph{check on a second point}. Work
the graph with VA $x=-1$, HA $y=-2$, through $(1,-1)$: $-1=\frac a2-2$ gives $a=2$, so
$y=\frac{2}{x+1}-2$, and $(0,0)$ confirms it. Flag the sign trap out loud: the asymptote is at $x=-1$
and the equation shows $x+1$, because $x-(-1)=x+1$. This is the likeliest SOL item on the standard; do
not let it get cut for time.

\textbf{6. Compare, contrast, and the disguise (A2.F.1a/e).} Fill the table against $y=x^2$ and
$y=|x|$ --- restricted domain, restricted range, asymptotes, no intercepts, two pieces, origin
symmetry. Then the payoff: $\frac{x-1}{x-3}=\frac{(x-3)+2}{x-3}=1+\frac{2}{x-3}$, which is the Section
4 anchor. Ask who recognizes the Warm-Up. Whenever top and bottom are both linear, the function is a
shifted, stretched $\frac1x$ --- and this is \emph{why} 4.0's degree rule is true in that case: the
leftover constant \emph{is} the ratio of the leading coefficients.
\end{multicols}
\end{skillbox}

\begin{skillbox}[Explicit Instruction: Graph It, Then Un-Graph It]{forestbg}
\begin{tabularx}{\linewidth}{@{}p{0.46\linewidth} X@{}}
\vspace{0pt}
\textbf{Equation $\rightarrow$ graph (A2.F.1c). Post these.}
\begin{enumerate}[leftmargin=1.3em, nosep, after=\vspace{2pt}]
  \item Set the denominator to zero: \textbf{VA} $x=h$.
  \item The number added outside is the \textbf{HA} $y=k$.
  \item Draw both as dashed lines. They cross at the \textbf{center} $(h,k)$.
  \item \textbf{Domain} $x\neq h$; \textbf{range} $y\neq k$.
  \item \textbf{Intercepts:} $y$-int is $g(0)$; $x$-int solves $\frac{a}{x-h}=-k$.
  \item Use $a$ to decide which pair of regions the branches occupy and how far out they sit.
\end{enumerate}

\vspace{3pt}
\textbf{Graph $\rightarrow$ equation (A2.F.1b).}
\begin{enumerate}[leftmargin=1.3em, nosep, after=\vspace{2pt}]
  \item Read the VA: $h$. \quad 2. Read the HA: $k$.
  \item[3.] Write $y=\frac{a}{x-h}+k$ with $h,k$ filled in.
  \item[4.] Substitute one marked point; solve for $a$.
  \item[5.] \textbf{Check on a second point.}
\end{enumerate}
&
\vspace{0pt}
\textbf{Worked, on the anchor.} \; $g(x)=\dfrac{2}{x-3}+1$
\[ x-3=0 \Rightarrow \text{VA } x=3; \qquad \text{HA } y=1; \qquad \text{center } (3,1) \]
\[ g(0)=\frac{2}{-3}+1=\frac13; \qquad \frac{2}{x-3}=-1 \Rightarrow x=1 \]
\centerline{Domain $x\neq3$ \quad Range $y\neq1$ \quad $a=2>1$: branches pushed out}

\vspace{4pt}
\textbf{And the same function, in disguise:}
\[ \frac{x-1}{x-3}=\frac{(x-3)+2}{x-3}=1+\frac{2}{x-3} \]
\vspace{-6pt}
\textbf{The two questions students must answer aloud:} \emph{which number moved which asymptote?} and
\emph{is $(h,k)$ on the graph?} (No --- it is the one point the graph is guaranteed to miss.)
\end{tabularx}
\end{skillbox}

\begin{skillbox}[Active Monitoring]{redbox}
Circulate during the guided practice and Tier work. Watch for and correct:
\begin{itemize}[leftmargin=1.2em, nosep]
  \item \textbf{the sign of $h$} --- $\frac{1}{x-3}$ read as a shift \emph{left}; the signature error, and worse here than in Units 1--2 because it puts the wall on the wrong side of the $y$-axis;
  \item \textbf{the horizontal asymptote read off the numerator} --- $a$ and $k$ are different letters doing different jobs, and $a$ moves nothing;
  \item \textbf{the two roles crossed} --- ``the $+3$ inside moves it up'' --- which ruins every problem after it;
  \item \textbf{``the parent goes through the origin''} --- it does not, and this is the fact students most resist;
  \item \textbf{range given as ``all reals''} --- domains get restricted by habit, ranges get forgotten; the HA is a \emph{range} exclusion;
  \item \textbf{decreasing on $(-\infty,\infty)$} --- 4.0's union rule, still in force, still costing points;
  \item \textbf{intercepts guessed off the picture} instead of computed from the equation;
  \item \textbf{$(h,k)$ plotted as a point of the graph} --- it is precisely the point the graph never reaches;
  \item \textbf{stalling on $a$} in a graph-to-equation item --- that student has the concept and is missing a substitution; one minute fixes it.
\end{itemize}
Cold-call prompts: ``Which asymptote did that number move, and how do you know?'' ``What is the domain,
and now what is the \emph{range}?'' ``Is the center on the graph?'' ``Your equation fits one point ---
now try the other one.'' ``$x-3$: left or right? Prove it with a number.'' ``Can this graph ever
\emph{touch} $y=1$?''
\end{skillbox}

\begin{skillbox}[Group Work \& Differentiation]{redbox}
\textbf{Move the Walls} activity (tiered; mirrors the notes).
\begin{multicols}{3}
\textbf{Tier R --- Remediate.} One transformation at a time. Four asymptote-naming items
($\frac{1}{x-5}$, $\frac1x-4$, $\frac{1}{x+6}$, then the full $\frac{2}{x-1}+3$ with center, domain,
and range), a \textbf{graph-matching} set of three pre-drawn hyperbolas against
$\frac1x+3$, $\frac{1}{x+2}-3$, $\frac{1}{x-2}$, and a from-memory recall of the parent's features. The
match is the item to watch: choosing the horizontally-shifted graph for $\frac1x+3$ means the two roles
are crossed, and that must be fixed before the student goes further.

\columnbreak
\textbf{Tier A --- Approaching.} A four-row table ($a,h,k$ / VA \& HA / domain \& range / $y$-intercept)
over $\frac{4}{x-1}$, $\frac1x-5$, $\frac{-2}{x+3}+1$, $\frac{6}{x+2}-2$ --- the second row is the
point, since $\frac1x-5$ has \emph{no} $y$-intercept because its VA is the $y$-axis. Then a
\textbf{graph-to-equation} item (VA $x=1$, HA $y=3$, through $(3,2)$, so $a=-2$) with a mandatory check
on the second point. Then an \textbf{error analysis}: a student reads $\frac{1}{x+4}-1$ as ``VA $x=4$,
HA $y=1$'' --- disprove the wall by evaluating at $x=4$ and getting $-\frac78$.

\columnbreak
\textbf{Tier E --- Extension.} \emph{Part 1:} unmask $\frac{3x+5}{x+1}=3+\frac{2}{x+1}$ and
$\frac{2x-7}{x-4}=2+\frac{1}{x-4}$, then confirm both horizontal asymptotes against 4.0's degree rule.
\emph{Part 2:} build one backwards --- VA $x=-3$, HA $y=4$, through $(-2,1)$ gives
$\frac{-3}{x+3}+4=\frac{4x+9}{x+3}$ --- plus the question of whether a graph with those asymptotes could
pass through $(-3,1)$ (it could not; $-3$ is not in the domain). \emph{Part 3:} \textbf{diluting the
brine} --- $20$ g of salt in $5$ L, add $x$ L of pure water, $C(x)=\frac{20}{5+x}$. The HA $y=0$ means
the salt never leaves; the VA $x=-5$ is real algebra and meaningless chemistry; and halving the
concentration costs $5$, then $10$, then $20$ more liters, which is diminishing returns made visible.
\end{multicols}
\end{skillbox}

\begin{skillbox}[Individual Work \& Assessment]{redbox}
\textbf{Exit Ticket} (independent, no notes): \textbf{(1)} read $g(x)=\frac{-3}{x+2}-1$ completely ---
$a,h,k$, center $(-2,-1)$, VA $x=-2$, HA $y=-1$, $y$-intercept $(0,-\frac52)$, domain $x\neq-2$, range
$y\neq-1$ --- and then answer whether the graph ever \emph{crosses} its horizontal asymptote (it does
not: that would need $\frac{-3}{x+2}=0$). \textbf{(2)} the \textbf{A2.F.1b item}: write the equation
from a graph with VA $x=2$, HA $y=-1$, through $(3,2)$ and $(5,0)$ --- answer $y=\frac{3}{x-2}-1$, with
the second point used as the check. \textbf{(3)} an \textbf{SOL-style multiple-choice} item: which
function has VA $x=-4$ and HA $y=3$? Each distractor breaks exactly one idea --- $\frac{2}{x-4}+3$
flips the sign of $h$, $\frac{2}{x+3}+4$ swaps the two jobs, and $\frac{3}{x+4}+2$ reads the horizontal
asymptote off the numerator. Collect and sort into ``got it / sign of $h$ / the two jobs crossed /
equation-from-a-graph.'' The sign pile is the fastest fix. The \emph{crossed-jobs} pile is the one to
re-teach before tomorrow, because 4.5 takes the asymptotes off the face of the equation entirely --- a
student who cannot say which letter moves which asymptote has nothing left to hold onto there.
\end{skillbox}

\begin{skillbox}[Reinforcement \& Extension]{goldbox}
\textbf{Homework:} six quick asymptote-and-move items covering each transformation singly and in
combination; three \textbf{full reads} ($\frac{4}{x-1}-2$, $\frac{-6}{x+3}+3$, $\frac{1}{x+5}$) with
center, asymptotes, domain, range, and both intercepts, closing on \emph{which one has no $x$-intercept
and why} (the third: $k=0$, so its HA is the $x$-axis); two \textbf{graph-to-equation} items, the
second of which deliberately omits the horizontal asymptote from the figure so students must read
$y=0$ off the flattening branches; four \textbf{disguises} ($\frac{x+5}{x+2}$, $\frac{4x-3}{x-1}$,
$\frac{2x}{x+4}$ --- the last has no visible constant to split off, and needs $2x=2(x+4)-8$) checked
against the degree rule; an \textbf{A2.F.1e table comparison} pairing $y=3x$ (constant differences)
against $y=\frac{12}{x}$ (constant \emph{products}) --- which is quietly an inverse variation and a seed
for 4.7; and a \textbf{free-throw model}, $P(x)=\frac{12+x}{20+x}=1-\frac{8}{x+20}$, where reaching
$90\%$ takes $60$ makes in a row, $95\%$ takes $140$, and $100\%$ takes forever --- a horizontal
asymptote priced in a currency students care about. \textbf{Extension:} three students describe
$\frac{-2}{x-5}+4$ --- one flips $h$, one reads the HA off the numerator, and the third gets both
asymptotes right and still misses the \emph{reflection}, which is the instructive case --- plus a
\emph{design} task (VA $x=4$, HA $y=-3$, $y$-intercept $(0,-2)$) and the question of why a point on the
horizontal asymptote could never have pinned down $a$. \textbf{Preview:} Lesson 4.5, \emph{Graphing
Rational Functions} --- the denominator factors into two pieces, the asymptotes stop being visible in
the equation, and the graph has to be assembled from a factored form, a degree comparison, the
intercepts, and a sign chart.
\end{skillbox}
\end{document}
